% year: תשע"ז
% type: מבחן
% semester: א
% moed: ב
% number: 4א
% points: 12
% categories: func-analysis
% source: exams/מבחנים/תשעז/מועד ב תשעז.pdf

\begin{question}
(12 נק') חקרו את הפונקציה $f(x)=x^2e^{-x}$ ושרטטו את הגרף שלה.
\underline{צריך למצוא}: תחום הגדרה, אסימפטוטות; תחומי עליה וירידה, נקודות קיצון, תחומי קמירות וקעירות.
\end{question}

\begin{solution}
\textbf{תחום הגדרה:} כל $\R$; הפונקציה רציפה וגזירה בכל נקודה, ו-$f(x)\ge0$ עם שוויון רק ב-$x=0$.

\textbf{אסימפטוטות:} אין אנכיות (רציפה על $\R$).
כאשר $x\to+\infty$: $\lim_{x\to\infty}\frac{x^2}{e^x}=0$ (לופיטל פעמיים), ולכן $y=0$ אסימפטוטה אופקית ב-$+\infty$.
כאשר $x\to-\infty$: $f(x)\to+\infty$ ו-$\frac{f(x)}{x}=xe^{-x}\to-\infty$, ולכן אין אסימפטוטה (אופקית או משופעת) ב-$-\infty$.

\textbf{עליה וירידה:}
\[
f'(x)=2xe^{-x}-x^2e^{-x}=x(2-x)e^{-x}.
\]
$f'<0$ ב-$(-\infty,0)$ וב-$(2,\infty)$ (יורדת), ו-$f'>0$ ב-$(0,2)$ (עולה).

\textbf{קיצון:} מינימום (מוחלט) ב-$(0,0)$; מקסימום מקומי ב-$\left(2,\frac4{e^2}\right)\approx(2,0.54)$ (לא מוחלט, כי $f\to\infty$ ב-$-\infty$).

\textbf{קמירות:}
\[
f''(x)=\big((2-2x)-(2x-x^2)\big)e^{-x}=(x^2-4x+2)e^{-x},
\]
שמתאפסת ב-$x=2\pm\sqrt2$. לכן $f$ קמורה ($\cup$) ב-$(-\infty,2-\sqrt2)$ וב-$(2+\sqrt2,\infty)$, וקעורה ($\cap$) ב-$(2-\sqrt2,2+\sqrt2)$. נקודות פיתול ב-$x=2\pm\sqrt2$ (בערך $x\approx0.59$ ו-$x\approx3.41$).

\textbf{סרטוט:} הגרף יורד מ-$+\infty$ (ב-$-\infty$) עד המינימום $(0,0)$ ומשיק לציר $x$ שם, עולה עד המקסימום $(2,4e^{-2})$, ואז יורד ושואף ל-$0$ מלמעלה ($y=0$ אסימפטוטה ב-$+\infty$).
\end{solution}
